\documentclass[12pt,letterpaper]{article}

% ── Packages ─────────────────────────────────────────────────────────────────
\usepackage{amsmath,amssymb,amsthm}
\usepackage{mathtools}
\usepackage{geometry}
\usepackage{hyperref}
\usepackage{booktabs}
\usepackage{enumitem}
\usepackage{microtype}
\usepackage[numbers,sort&compress]{natbib}

\geometry{margin=1in}

% ── Theorem environments ──────────────────────────────────────────────────────
\theoremstyle{plain}
\newtheorem{theorem}{Theorem}[section]
\newtheorem{lemma}[theorem]{Lemma}
\newtheorem{proposition}[theorem]{Proposition}
\theoremstyle{definition}
\newtheorem{definition}[theorem]{Definition}
\newtheorem{assumption}[theorem]{Assumption}
\newtheorem{prediction}[theorem]{Prediction}
\theoremstyle{remark}
\newtheorem{remark}[theorem]{Remark}

% ── Custom commands ───────────────────────────────────────────────────────────
\newcommand{\tr}{\operatorname{Tr}}
\newcommand{\kB}{k_{\mathrm{B}}}
\newcommand{\lnTwo}{\ln 2}
\newcommand{\Qmin}{Q_{\min}}
\newcommand{\Smax}{S_{\max}}

% ── Metadata ─────────────────────────────────────────────────────────────────
\title{\textbf{Landauer's Principle in the Emergent Thermodynamic Information Framework}\\[4pt]
\large Operational Foundations, Rigorous Limits, and Testable Predictions}

\author{Kevin Monette\\[2pt]
\small Independent Researcher (AI-assisted research)\\
\small \texttt{kevin.monette@research.org}}

\date{February 9, 2026 \quad (repository version March 2026)}

% ─────────────────────────────────────────────────────────────────────────────
\begin{document}
\maketitle

\begin{abstract}
Landauer's principle --- that erasing one bit of information at temperature $T$ dissipates at
least $\Qmin = \kB T \lnTwo$ --- is commonly mischaracterized as either a fundamental law of
nature or a speculative claim.  We show it is neither: it is a precise thermodynamic
\emph{cost of agency}, a constraint on how physical systems implement logically irreversible
operations.

Within the Emergent Thermodynamic Information (ETI) framework, we derive Landauer's principle
from five explicit assumptions~(A1--A5), state four lemmas~(L1--L5) as rigorous consequences,
and formulate three falsifiable predictions~(P1--P3).  We further show that (i)~vacuum
fluctuations do not violate the principle, (ii)~reversible computation defers but cannot
eliminate dissipation, and (iii)~the universe as a whole cannot violate the principle because it
has no external entropy sink.

The framework is non-mystical, observer-independent, and grounded in established quantum
thermodynamics.  The ETI treatment of information as emergent from physical correlations and
constraints --- never as a fundamental substance --- clarifies longstanding confusions about the
relationship between information, entropy, and energy.
\end{abstract}

\tableofcontents
\bigskip

% ─────────────────────────────────────────────────────────────────────────────
\section{Foundational Assumptions}
\label{sec:assumptions}

All results derive from the following five assumptions, stated explicitly so that any claim can
be traced to its premises.

\begin{assumption}[Causal Closure]\label{ass:A1}
The universe $\mathcal{U}$ is a closed, causally connected system.  No external entropy sinks
exist outside $\mathcal{U}$.
\end{assumption}

\begin{assumption}[Microdynamics]\label{ass:A2}
Closed systems evolve unitarily under $U(t)$ on Hilbert space $\mathcal{H}$.  Open subsystems
evolve via completely positive trace-preserving (CPTP) maps $\mathcal{E}$ on density operators.
\end{assumption}

\begin{assumption}[Thermodynamics as Effective]\label{ass:A3}
The thermodynamic entropy $S(\rho) = -\kB\tr(\rho\ln\rho)$ is a coarse-grained statistical
description of a physical state relative to a chosen partitioning.  It is emergent, not
fundamental.
\end{assumption}

\begin{assumption}[Physical Memory]\label{ass:A4}
Logical information (bits, qubits) is always instantiated in physical substrates with stability
requirements: memory states must be distinguishable, persistent on relevant timescales, and not
spontaneously decohered faster than the operation rate.
\end{assumption}

\begin{assumption}[Finite Resources]\label{ass:A5}
Any physical agent operates under finite memory, finite energy and cooling capacity, and finite
control precision.  Indefinite information storage without resource expenditure is impossible
under Assumptions~\ref{ass:A1}--\ref{ass:A4}.
\end{assumption}

% ─────────────────────────────────────────────────────────────────────────────
\section{Definitions}
\label{sec:definitions}

\subsection{Logical vs.\ Physical Operations}

Let a memory register have logical state space $\mathcal{M} = \{0,1\}^{n}$, implemented via
physical Hilbert space $\mathcal{H}$.

\begin{definition}[Logically Irreversible Operation]
A map $f\colon\mathcal{M}\to\mathcal{M}$ is \emph{logically irreversible} if it is
many-to-one:
\begin{equation}
  \exists\; m \neq m' \in \mathcal{M} \;\text{ such that }\; f(m) = f(m').
  \label{eq:irreversible}
\end{equation}
The canonical example is the reset operation $f(0) = f(1) = 0$.
\end{definition}

\begin{definition}[Logically Reversible Operation]
A bijection on $\mathcal{M}$, implementable by a unitary $U$ on $\mathcal{H}$ acting as a
permutation on the physical states corresponding to $\mathcal{M}$.
\end{definition}

\begin{remark}
Physical evolution of a closed system is always reversible (unitary).  Thermodynamic cost
arises from the \emph{choice to lose information} about prior states, not from physics per se.
\end{remark}

\subsection{Entropy and Negentropy}

\begin{definition}[Thermodynamic Entropy]
\begin{equation}
  S(\rho) = -\kB\,\tr(\rho\ln\rho).
  \label{eq:vN-entropy}
\end{equation}
\end{definition}

\begin{definition}[Negentropy]
\begin{equation}
  N(\rho) = S(\rho_{\max}) - S(\rho),
  \label{eq:negentropy}
\end{equation}
where $\rho_{\max}$ is the maximally mixed state on the same support.  Negentropy measures
local deviation from maximum disorder.
\end{definition}

\begin{remark}
Negentropy is not conserved; it is not Shannon information; and in the ETI framework,
``information'' is not a fundamental substance but an emergent descriptor of physical
correlations and constraints.
\end{remark}

% ─────────────────────────────────────────────────────────────────────────────
\section{Landauer's Principle: Operational Statement}
\label{sec:landauer}

\subsection{Standard Formulation}

\begin{theorem}[Landauer 1961, rigorous form Bennett 2003]\label{thm:landauer}
Resetting a single bit of information stored in a physical memory in thermal equilibrium at
temperature $T$ requires dissipation of at least
\begin{equation}
  Q \geq \kB T \lnTwo
  \label{eq:landauer-bound}
\end{equation}
into an effective thermal reservoir, under the following conditions:
\begin{enumerate}[label=(\roman*)]
  \item the memory is in thermal equilibrium with a bath at temperature $T$;
  \item memory states are energetically degenerate, distinguishable, and stable
        (Assumption~\ref{ass:A4});
  \item the reset is logically irreversible (a many-to-one map);
  \item the reset succeeds with high probability.
\end{enumerate}
\end{theorem}

\subsection{The Correct Interpretation}

Landauer's principle is a constraint on the thermodynamic cost of implementing logically
irreversible operations on physical substrates.  It states precisely: if information about the
prior state is lost in a logically irreversible way, entropy must be exported to the environment
at a minimum rate of $\kB\lnTwo$ per bit erased at temperature $T$.

It does \emph{not} state: ``information cannot be erased'' or ``information is physical in a
mystical sense.'' The cost arises because erasure is a \emph{physical process} that must
redistribute entropy; it is not because information is a substance.

% ─────────────────────────────────────────────────────────────────────────────
\section{Reversible Computation and the Persistence of Dissipation}
\label{sec:reversible}

\subsection{Ideal Reversible Gates: No Immediate Cost}

A unitary gate $U$ acting on a pure state $|\psi\rangle$ produces a pure state.  No entropy is
generated:
\begin{equation}
  S\!\left(U\rho U^{\dagger}\right) = S(\rho)
  \quad \text{for any unitary }U.
  \label{eq:unitary-entropy}
\end{equation}
Reversible gates do not require dissipation in the logical transformation.  They defer
dissipation but cannot eliminate it.

\subsection{Why Sustained Computing Dissipates}

Even with all-reversible gates, three mechanisms force entropy export under
Assumption~\ref{ass:A5}:

\paragraph{Error correction.}
Quantum error correction requires syndrome extraction (measurement) and ancilla reset.  Each
reset incurs a Landauer cost.  In the surface code, each syndrome measurement requires resetting
ancilla qubits, costing at least $\kB T\lnTwo$ per bit per cycle.

\paragraph{Memory recycling.}
Any agent with finite memory must eventually overwrite registers.  Each overwrite is a logically
irreversible map: a Landauer operation.

\paragraph{Control and refrigeration.}
Maintaining quantum coherence requires active cooling and control-field generation --- both
thermodynamically irreversible processes generating waste heat in control infrastructure.

% ─────────────────────────────────────────────────────────────────────────────
\section{Lemmas: Rigorous Consequences}
\label{sec:lemmas}

\begin{lemma}[No External Entropy Sink]\label{lem:L1}
Under Assumption~\ref{ass:A1}, any entropy sink exchanging energy or information with
$\mathcal{U}$ is part of $\mathcal{U}$.  No external reservoir exists.
\end{lemma}

\begin{proof}
Any system exchanging energy or information with $\mathcal{U}$ is causally connected to
$\mathcal{U}$ and therefore, by Assumption~\ref{ass:A1}, part of $\mathcal{U}$.
\end{proof}

\begin{lemma}[Landauer Cost for Irreversible Reset]\label{lem:L2}
Under Assumptions~\ref{ass:A1}--\ref{ass:A4}, any implemented many-to-one reset of a stable
memory in thermal equilibrium at temperature $T$ incurs entropy export of at least
$\kB\lnTwo$ per bit to the environment.
\end{lemma}

\begin{proof}[Proof sketch]
The reset maps two distinct microstates to one.  Under unitarity (Assumption~\ref{ass:A2}),
degrees of freedom must be transferred from the memory to the environment to preserve total
state-space dimension.  This expansion of the environment's accessible microstates constitutes
an entropy increase of at least $\kB\lnTwo$ per bit, by the argument of Penrose and Bennett
\cite{Bennett2003}.
\end{proof}

\begin{lemma}[Reversible Computation Defers Dissipation]\label{lem:L3}
Under Assumption~\ref{ass:A2}, unitary gates do not require dissipation in the logical
transformation.  Under Assumption~\ref{ass:A5}, dissipation is inevitable for sustained
finite-resource computation via error correction, memory recycling, or control overhead.
\end{lemma}

\begin{lemma}[Sustained Computing Requires Entropy Export]\label{lem:L4}
Under Assumptions~\ref{ass:A2} and~\ref{ass:A5}, any computing process on $n$ qubits with
error rate $\Gamma_{\mathrm{err}}$ per qubit running for time $t$ must export total entropy at
least:
\begin{equation}
  \Delta S_{\mathrm{env}}
  \;\geq\;
  \Gamma_{\mathrm{err}}\cdot n \cdot t \cdot \kB\lnTwo.
  \label{eq:entropy-export}
\end{equation}
\end{lemma}

\begin{lemma}[Vacuum Fluctuations Are Not Free Thermodynamic Fuel]\label{lem:L5}
Vacuum fluctuations do not violate Landauer's principle, because they are not logical
operations.  They do not reset, overwrite, or record information in a way requiring a
many-to-one mapping on logical states.  Any thermodynamic relevance requires coupling to an
apparatus that measures, stores, and eventually resets --- incurring Landauer cost at those
stages.
\end{lemma}

\begin{remark}
Lemma~\ref{lem:L5} rules out Casimir-based or zero-point-energy-based ``free entropy''
proposals.  Any scheme claiming sustained net work from vacuum fluctuations must identify where
the Landauer cost of the measurement--storage--reset cycle is paid.
\end{remark}

% ─────────────────────────────────────────────────────────────────────────────
\section{Predictions: Falsifiable Claims}
\label{sec:predictions}

\begin{prediction}[Scaling of Coherent Computation Dissipation]\label{pred:P1}
As quantum computers scale to larger systems and longer computation times, the total system
entropy export (cooling load plus error-correction overhead) increases with system size and
computation length.  No asymptotic limit at which a quantum computer of size $n$ qubits
running for time $t$ requires zero entropy export exists under Assumption~\ref{ass:A5}.

\medskip\noindent
\textit{Test:} Measure total cooling power required for superconducting quantum computers of
sizes $n \in \{10, 50, 100, 1000\}$ qubits at fixed error rate, normalized per logical gate
operation.
\end{prediction}

\begin{prediction}[Vacuum Work Extraction]\label{pred:P2}
Any proposal claiming to extract net useful work from vacuum fluctuations indefinitely must
specify: (i)~the initial state of the extraction apparatus, and (ii)~where the Landauer cost
of the measurement--storage--reset cycle is paid.  Under Lemma~\ref{lem:L5}, no such scheme
produces sustained net output when all entropy accounts are closed.

\medskip\noindent
\textit{Test:} Examine any specific vacuum energy extraction device.  Prediction: the entropy
export is either (a)~hidden in an ignored component, (b)~deferred to a periodic
reinitialization, or (c)~implicitly violating Assumption~\ref{ass:A1}.
\end{prediction}

\begin{prediction}[Sub-Landauer Erasure Claims]\label{pred:P3}
If any experiment reports bit erasure with energy cost below $\kB T\lnTwo$ per bit, at least
one of the following must hold:
\begin{enumerate}[label=(\roman*)]
  \item The temperature $T$ is defined non-standardly;
  \item The erasure success probability is not unity;
  \item Entropy is exported to a subsystem excluded from the energy budget;
  \item The system starts outside thermal equilibrium, consuming non-equilibrium resources.
\end{enumerate}
\end{prediction}

% ─────────────────────────────────────────────────────────────────────────────
\section{Integration with the Entropic Gravity Program}
\label{sec:gravity}

Within the ETI framework, measurement-induced constraint-fixing (formalized in the companion
paper~\cite{MonetteB2026}) is a Landauer operation: it irreversibly fixes a constraint and
exports entropy to the environment.  This entropy export creates a local negentropy gradient
$\nabla N$.

The companion entropic gravity paper~\cite{MonetteA2026} proposes that negentropy gradients
source attractive gravitational curvature through an extension of Jacobson's thermodynamic
gravity~\cite{Jacobson1995}.  If correct, every measurement --- every thermodynamically
irreversible constraint-fixing event --- produces a minute gravitational imprint.

This is the unifying thread of the research program: information processing and spacetime
geometry are the same phenomenon viewed at different scales.

% ─────────────────────────────────────────────────────────────────────────────
\section{Conclusion}
\label{sec:conclusion}

Landauer's principle is not a fundamental law of nature in the sense of Newton's laws.  It is
a precise, operational consequence of implementing logically irreversible operations on physical
substrates in a universe with no external entropy sink.

\medskip

\begin{center}
\textit{Landauer's principle is not a law --- it is a cost of agency.}
\end{center}

\medskip

Whenever a physical agent exercises control in a logically irreversible way, thermodynamic
entropy must be redistributed to the environment at the minimum rate set by
Theorem~\ref{thm:landauer}.  This is not mystical.  It is the cleanest possible statement of
the relationship between physical processes and information management in a causally closed
universe.

% ─────────────────────────────────────────────────────────────────────────────
\section*{Acknowledgments}
The author acknowledges AI research partners for assistance with literature synthesis and
critical review.  All theoretical content originates from the author's independent research
program.

% ─────────────────────────────────────────────────────────────────────────────
\bibliographystyle{unsrtnat}
\begin{thebibliography}{99}

\bibitem{Landauer1961}
R.~Landauer,
``Irreversibility and heat generation in the computing process,''
\textit{IBM J.\ Res.\ Dev.}\ \textbf{5}, 183 (1961).
\href{https://doi.org/10.1147/rd.53.0183}{doi:10.1147/rd.53.0183}

\bibitem{Bennett2003}
C.~H.~Bennett,
``Notes on Landauer's principle, reversible computation, and Maxwell's demon,''
\textit{Stud.\ Hist.\ Phil.\ Mod.\ Phys.}\ \textbf{34}, 501 (2003).
\href{https://doi.org/10.1016/S1355-2198(03)00039-X}{doi:10.1016/S1355-2198(03)00039-X}

\bibitem{Jacobson1995}
T.~Jacobson,
``Thermodynamics of spacetime: The Einstein equation of state,''
\textit{Phys.\ Rev.\ Lett.}\ \textbf{75}, 1260 (1995).
\href{https://doi.org/10.1103/PhysRevLett.75.1260}{doi:10.1103/PhysRevLett.75.1260}

\bibitem{Parrondo2015}
J.~M.~R.~Parrondo, J.~M.~Horowitz, and T.~Sagawa,
``Thermodynamics of information,''
\textit{Nature Phys.}\ \textbf{11}, 131 (2015).
\href{https://doi.org/10.1038/nphys3230}{doi:10.1038/nphys3230}

\bibitem{Sagawa2009}
T.~Sagawa and M.~Ueda,
``Minimal energy cost for thermodynamic information processing,''
\textit{Phys.\ Rev.\ Lett.}\ \textbf{102}, 250602 (2009).
\href{https://doi.org/10.1103/PhysRevLett.102.250602}{doi:10.1103/PhysRevLett.102.250602}

\bibitem{MonetteA2026}
K.~Monette,
``Gravitational coupling to entanglement entropy density,''
Independent Research Program, March 2026.

\bibitem{MonetteB2026}
K.~Monette,
``Constraint manifolds and the limits of quantum observability,''
Independent Research Program, March 2026.

\end{thebibliography}

\end{document}
